\section{Change log and study history}
\label{v5-sec:history}\label{sec:v5-history}
This appendix begins with the history of the studies used to prepare the
version 5.0.0 unblinding-request draft. The main text is a curated account
of the current proposed procedure, while this chronology explains how the
support choices, background treatment and interpretation changed.
A later report does not automatically supersede every earlier result:
the underlying data fraction, fit policy, numerical qualification and
probability definition determine which comparison is valid.

\subsection{Historical validation and the full-note foundation}
The v4 analysis established the sideband GP interpolation, signal
extraction and bounded asymptotic $CL_s$ workflow. Subsequent v4.6--v4.8
studies examined trigger categories, injected-signal response, guard bands,
residual structure and smooth alternative backgrounds. Their plots remain
useful historical controls. They are not measurements of observed-data
bias or global significance, and tests with different truth constructions
or injected-yield definitions are not combined into a single ensemble.

The August 2026 v4.9 and v4.9.1 threshold studies established that 2021
signal recovery depends on the lower training edge and on the generating
shape. The v4.9.5 support study, recorded on 20 August, compared six
lower edges. No choice passed its original uniform mean-pull criterion.
A documented amendment selected the 36~MeV edge under a practical
criterion and confirmed it with independent backgrounds. The remaining
55~MeV offsets are retained in Appendix~\ref{v5-sec:calibration}.
The older 40--300~MeV 2021 support is historical context; the current
2021 card uses the documented 36--300~MeV support.

The v4.9.7 study of 2 September attempted the 2016 support qualification
needed for subsequent combined extraction. Its initial support scan failed
the declared gate and stopped; its planned continuation and combined
validation were not run. A later, separate observed-state replay retained
an explicit 2016 numerical exception: 87 of 142 states were resolved,
with 55 unresolved. These later qualifications take precedence over any
historical statement that the validation alone established readiness for
full-data unblinding. Version 5 retains the technical description of the
detector inputs, signal normalization and fitting procedure from that
full-note lineage, while stating the remaining review conditions explicitly.

The Harvard final-combinations writing sample, also prepared on
2 September, supplied the compact presentation of the available observed
results and polished explanatory prose. Its inputs were full 2015,
full 2016 and the released 2021 10\% sample. This draft expands those
sections with the subsequent statistical studies and separates the
proposal to unblind from a claim that all calibration gates have passed.
It does not open additional 2021 events.

\subsection{Observed combinations, tails and resolution: v4.9.12}
The 4--5 September v4.9.12 sequence consolidated observed combinations,
expected-limit reference products, targeted statistical-tail sampling,
and mass-resolution variations. Version 5 uses its observed individual
and combined coupling limits, supplements the principal figure with
asymptotic local probabilities, and retains the requested resolution
comparisons as response studies. Figures from different releases are
identified by their numerical source rather than assumed to be equivalent
because their appearance is similar.

The targeted-tail study refined finite-simulation probabilities where
possible. A zero exceedance count supplies a sampling bound and does
not measure a zero probability; an unresolved endpoint remains
explicitly unresolved. Resolution variations and selected candidate
extractions assess how the result responds to a changed template.
They do not independently measure a resolution error or supply an
additional significance gain.

The related 2021 peak--dip and reverse-injection studies showed how a
positive localized feature can enter neighboring sidebands and induce
negative or positive fitted responses as the mask moves. These studies
motivate the correlation and extraction displays in the results section.
Their generating smooth shapes and selected masses make them conditional
response diagnostics.

\subsection{Background profiling and empirical calibration: v4.9.13}
The 5 September background-profiling study compared the nominal
correlated Gaussian profile, a positive log-GP profile and a fixed GP
mean. Known-background controls separated implementation behavior from
uncertainty propagation; injections with fluctuating background and
retrained sidebands demonstrated the limitations of the fixed approach
and the remaining 71~MeV failure of the profile.
Version 5 brings the controlled limit comparison and injected-bias figure
into the methodology discussion, including the physical meaning of the
two- and five-reference-error injections.

The accompanying calibration release evaluated 456 dataset/mass
coordinates and independently checked exclusion and local-rejection
rates. Its finite-set background envelope removed the adjusted failures
within the declared validation family, while exposing large sensitivity
to the 2016 stress construction. Appendix~\ref{v5-sec:calibration}
reorganizes the full calibration discussion around a flowchart,
background definitions, observed changes, interpolation offsets and
validation. Its conclusion supports profiling but does not certify
uncorrected asymptotic probabilities across every tested background.

\subsection{Understanding calibration and extending global tests: v4.9.14--v4.9.15}
The 6 September v4.9.14 interpretation report explained why calibration
can turn an unusually tight limit into a broad feature, and why a large
background-dependent offset differs from a wider statistical error.
It also introduced the 2015 significance-GP pilot, using ten complete
initial scans, 1,000 independent validation scans and 200,000 correlated
fields. Version 5 adopts its distinction between the background GP and
the GP describing correlations across the significance scan.

The v4.9.15 extension repeated this calculation for full 2016 and
released 2021 10\% data. It documented the coherent background
construction, central-distribution checks and unresolved rare tails.
These results motivate the method in Sec.~\ref{v5-sec:global-method},
while the source-fit and convergence limitations of the 2016 stress spectrum prevent
its very small formal local probabilities from becoming particle claims.
The 2015 and 2021 finite global-tail estimates remain conditional
method-development results.

\subsection{Union scan, probability audit and candidate studies: v4.9.16}
The 6 September combined-global study extended the actual shared-coupling
likelihood across the complete 19--250~MeV union. Dataset membership changes
with mass, and the new coherent calculation preserves correlations across
those boundaries. It supplies 232 observed mass points and separates
raw-local ordering from stress-offset-adjusted ordering. Its union limits
remain pointwise asymptotic limits; the 41-point earlier calibration
cannot be silently promoted to a calibration of the entire union.

The probability-audit and echo report checked the displayed extraction
likelihoods against the saved scan. It showed that the apparently extreme
76~MeV stress-reference result accompanies a small observed positive
local fit. Version 5 uses that distinction to label local significance
and conditional background-reference probabilities consistently. Its
deterministic signal-injection replay also demonstrates that a real
injected feature can induce nearby oscillations through the fitting
procedure, without establishing that the observed features are particles.

The extraction-presentation report provided the candidate count and
residual displays used in the results section. The deficit extension
reused the same coherent experiments and Gaussian fields in the negative
direction. It is retained as a background diagnostic, with its scan plot
shown beside the extraction discussion and its detailed interpretation
in the appendix. It adds no independent toy ensemble or direction-selection
correction.

The candidate-removal and traditional-search report replaced selected
regions with conditional expectations and compared fits to the original
counts with conventional local backgrounds. These interventions changed
the peak--dip pattern but did not eliminate all remote oscillation.
The conventional fits were strongly dependent on window and polynomial
choice. Version 5 keeps their extraction displays as model comparisons
and their detailed study in the appendix. They are not independent
confirmation of candidates selected from the same data.

Finally, the 2015 low-mass side study searched 15--20~MeV with local
12--28~MeV GP support and nearby support and polynomial controls.
It checked the 17.25~MeV fluctuation and kernel-ceiling dependence,
using small conditional toy ensembles. The study is placed at the end
of the appendix, outside the qualified production range: the accepted
signal shape and resolution below 19~MeV still require validation.

\subsection{What version 5 changes, and the remaining review decision}
The new note assembles these studies into one account: detector and
statistical method, observed individual and combined results, candidate
and resolution diagnostics, then organized technical appendices.
The proposed primary fit profiles correlated background uncertainty;
the calibration and significance-GP work remain separately identified
alternatives and validation studies. Historical readiness language is
replaced by the concrete conditions in Sec.~\ref{sec:v5-review}.
The statistical-tail gaps and background qualifications remain visible.

The 8 September editing pass curates existing results and improves their
presentation. It does not supply new pseudoexperiments, a new global
particle significance, or an authorization to inspect the remaining
2021 data. The appendix contents that follow provide a guide to the
supporting studies, while the accompanying source ledger records the
release directories and artifacts used for this draft.
